Phase I's methodology is the paper's main methodological contribution:
a way of running negative experiments in which the negative is
attributable. Five practices recur in every era.

\paragraph{Frozen protocols.}
A protocol document (configuration, arms, measurement definitions,
falsifier inequalities, verdict rule) is committed before any
implementation. Amendments are separate user-approved documents. The
protocol text, not the result, governs the verdict; the CLLA bundle
protocol \cite{cllaprot} is the canonical example, with falsifiers
F1--F7 and a binary verdict rule (\S\ref{sec:clla}).

\paragraph{Identity gates.}
Every mechanism is implemented behind a configuration flag whose
default is off. Before any flag-on run is interpreted, a flag-off run
must reproduce the committed baseline byte-for-byte: the event-stream
FNV-1a hash and every snapshot frame. The CLLA-era identity anchor is
the same-seed alternated baseline (152{,}254-row event stream, hash
\texttt{9647ea8a0ca4dbd2}), reproduced exactly by the unbounded,
bounded, rule, reserve, first-exposure, and source-correction
implementations \cite{cllarecord,bbounded,rule,reserve,fe,fec}. The
recruitment-gain implementation additionally re-verified flag-off
byte-identity against the corrected-fe baseline before its matrix
\cite{rg8}. This gate makes ``the flag did nothing'' a measured
condition rather than a claim.

\paragraph{Deterministic replay and read-only audits.}
Every run is preserved under \texttt{runs/<id>/} with telemetry chunks,
snapshots (1 s cadence), metrics, and report. Audits --- the
information-bottleneck audit \cite{synmem}, the V2.1 information audit
\cite{infoaudit}, the weight-growth, drive, and recruitment audits
\cite{wgrowth,drive} --- are read-only instruments over committed
artifacts: they never re-run the organism. This lets later epochs
re-examine earlier epochs' claims at full resolution.

\paragraph{Falsifiers and preservation of failure.}
Experiments pre-register the inequality whose failure destroys the
hypothesis (e.g., CLLA F1--F7; allocation S1 second-block/first-block
mass ratio \(\ge 0.5\)). Failures are preserved with their runs,
including aborts: 6 of 12 unbounded-CLLA runs aborted on the runaway
detector, and the record reports them rather than re-running the
survivors \cite{cllarecord}. Abort-loaded verdicts follow the frozen
rule (``uninformative $\Rightarrow$ FAIL'').

\paragraph{Exploratory versus formal experiments.}
X-series and SDE-series probes run without E-numbers and cannot carry
falsifier verdicts; formal E-numbered experiments and CLLA-era
registered experiments carry the verdict language. Promotion requires
an interesting, discriminating, reproducible result \cite{overview}.
The recruitment-gain experiment (§\ref{sec:rg}) is a registered
experiment without an E-number, closed under its own outcome rules.

\paragraph{Record corrections preserved.}
Where the record itself was defective, the correction is part of the
record, not a silent edit. Two are scientifically relevant and appear in
this paper: the V2.1 Stage-C validation table was committed without
executing any Stage-C runs; the information audit discovered that the
referenced runs do not exist, retracted the verdict as ungrounded, and
added an integrity regression test preventing phantom run references
\cite{infoaudit,overview}. And the CLLA capability protocol's F2/F6
measurement windows (presentations 21--40 of each pattern) are empty by
construction of the interleaved curriculum (20 presentations per
pattern); the capability record reports F2/F6 as unmeasurable under the
frozen definition rather than substituting a window, and reports the
closest available window as a non-decisive diagnostic \cite{capability}.
A third correction --- the first-exposure allocator source vacuity ---
is handled in \S\ref{sec:allocation} because it changed the mechanism
under test.